% This conditional wrapper lets the section compile independently for preview.
% When pasted into the full acmart report, the document class is already loaded
% and the wrapper has no effect.
\makeatletter
\newif\ifmpthreestandalone
\@ifclassloaded{acmart}{\mpthreestandalonefalse}{%
  \mpthreestandalonetrue
  \documentclass[sigconf,nonacm]{acmart}
  \settopmatter{printacmref=false}
  \renewcommand\footnotetextcopyrightpermission[1]{}
  \pagestyle{plain}
  \setlength{\emergencystretch}{1em}
  \raggedbottom
  \title{CS598JBR MP3 Report Section}
  \author{Team 16}
  \begin{document}
  \maketitle
}
\makeatother

\section{MP3}

We used DeepSeekCoder-6.7B-instruct with greedy decoding on the same 20
program IDs used in MP1 and MP2. For both tasks, we compared a vanilla prompt
with a crafted prompt and generated one deterministic response per program.
The model was loaded with 4-bit NF4 quantization, and generation used a limit
of 1,536 new tokens with sampling disabled. Thus, each comparison changes the
prompt configuration rather than the decoding strategy. For either task, the
success rate is the number of entries with \texttt{is\_correct=True} divided
by 20.

\subsection{MP3-Task1: Python-to-Java Code Translation}

\subsubsection{Evaluation and Success Rate}\hfill\break
For each program, we supplied the HumanEval-X Python prompt and canonical
solution and asked the model to produce a Java \texttt{class Solution}. We
extracted the code between \texttt{[Java Start]} and \texttt{[Java End]},
performed only structural normalization needed by the HumanEval-X harness,
compiled it with \texttt{javac}, and ran the supplied Java tests. A translation
was correct only if it compiled and every test passed.

The normalization step removed package and generated import declarations,
supplied a fixed set of standard Java imports, and retained one balanced class.
If the model used another wrapper name, the harness renamed it to
\texttt{Solution}; if it returned only a method, the harness placed that method
inside \texttt{class Solution}. These operations are permitted interface
repairs and do not change the generated method body or algorithm. We recorded
compilation and execution failures separately so that syntactic/interface
problems would not be confused with behavioral errors.

The vanilla prompt translated 9 of 20 programs successfully (45\%), while the
crafted prompt translated 14 of 20 (70\%). This is an absolute improvement of
25 percentage points (pp), within the expected 20--50 pp range. Table~\ref{tab:mp3t1pairwise}
shows the pairwise outcomes and distinguishes failures that did not reach test
execution from failures observed after successful compilation.

\begin{center}
\begin{minipage}{\columnwidth}\centering
\scriptsize
\setlength{\tabcolsep}{2pt}
\begin{tabular}{@{}lcc@{}}
\toprule
\texttt{problem\_ID} & Vanilla & Crafted \\
\midrule
HumanEval/60  & True                      & True                      \\
HumanEval/24  & True                      & True                      \\
HumanEval/67  & False (test failure)      & False (compilation issue) \\
HumanEval/145 & False (test failure)      & False (test failure)      \\
HumanEval/129 & False (compilation issue) & False (test failure)      \\
HumanEval/77  & False (compilation issue) & \textbf{True}             \\
HumanEval/79  & True                      & True                      \\
HumanEval/20  & False (compilation issue) & \textbf{True}             \\
HumanEval/62  & False (compilation issue) & \textbf{True}             \\
HumanEval/102 & True                      & True                      \\
HumanEval/8   & False (test failure)      & False (test failure)      \\
HumanEval/81  & True                      & True                      \\
HumanEval/140 & True                      & True                      \\
HumanEval/161 & False (compilation issue) & \textbf{True}             \\
HumanEval/136 & False (compilation issue) & False (test failure)      \\
HumanEval/75  & True                      & False (compilation issue) \\
HumanEval/148 & False (test failure)      & \textbf{True}             \\
HumanEval/118 & True                      & True                      \\
HumanEval/126 & True                      & True                      \\
HumanEval/120 & False (compilation issue) & \textbf{True}             \\
\midrule
Correct       & 9/20                      & 14/20                     \\
Success rate  & 45\%                      & \textbf{70\%}             \\
\bottomrule
\end{tabular}
\par\smallskip
\refstepcounter{table}\label{tab:mp3t1pairwise}\textbf{Table \thetable:}
Pairwise translation outcomes. Bold marks a False$\rightarrow$True change.
The failure labels distinguish compilation issues from test failures.
\end{minipage}
\end{center}

The pairwise view gives more information than the aggregate percentage. Eight
programs passed under both prompts, six changed from failure to success, five
failed under both prompts, and one regressed. Therefore, the net gain of five
programs is the result of six improvements offset by the HumanEval/75
regression; it is not caused by a different sample of programs.

\subsubsection{Prompt Changes}\hfill\break
The vanilla prompt followed the handout's direct-translation example, with one
operational addition asking for a complete \texttt{class Solution}; the Python
program supplied to both settings included its docstring. It required the two
Java markers but supplied no target declaration. The crafted prompt
added the complete Java declaration from HumanEval-X, including the required
method name, parameter types, and return type. It instructed the model to
preserve Python boundary cases and early returns; explicitly check collection,
string, indexing, slicing, numeric, division, modulo, sorting, and mutation
semantics; and avoid a \texttt{main} method, tests, a package declaration, or a
public top-level class. It also required a self-contained answer with standard
library imports when needed. These additions gave the model both the target
interface and a checklist of common cross-language semantic mismatches.

The declaration was especially useful because equivalent Python and Java
solutions need not use mechanically identical interfaces. Python functions can
rely on dynamic types, slicing, and implicit numeric behavior, whereas the Java
harness expects an exact method name and statically typed parameters. Supplying
the declaration separated the task of choosing the required interface from the
task of translating the algorithm. The checklist then directed attention to
semantic differences that a correct declaration alone cannot resolve.

\subsubsection{Translation Successes and Failures}\hfill\break
Compilation was the most prevalent vanilla failure: 7 translations failed to
compile, compared with 4 that compiled but failed tests. With the crafted
prompt, compilation failures fell from 7 to 2, while test failures remained at
4. Five of the six improvements (HumanEval/20, /62, /77, /120, and /161)
converted a vanilla compilation failure into a passing translation; the sixth
(HumanEval/148) corrected a behavioral failure. This shows that supplying the
exact Java declaration and output constraints primarily improved structural
compatibility.

The remaining failures show that compilation is not sufficient for semantic
equivalence. HumanEval/129 and /136 progressed from compilation errors to
executable translations but still failed tests, while HumanEval/8 and /145
failed tests under both prompts. These programs exercise state or index
tracking, sign-sensitive behavior, aggregate values, or custom ordering, all of
which require more than a surface-level syntax translation. Prompt crafting was
also not monotonic: HumanEval/67 changed from a test failure to a compilation
failure, and HumanEval/75 regressed from a pass to a compilation failure. Thus,
the persistent pattern is that the target declaration removes many interface
errors, but detailed control-flow and language-semantics errors remain.

Inspection of the generated Java made these categories more concrete.
HumanEval/67 did not translate the intended digit aggregation: the vanilla
response tried to parse the words around \texttt{apples} and \texttt{oranges},
while the crafted response contained a malformed array-index expression and
did not compile. HumanEval/75 was a truncation failure: the crafted response
ended in the middle of a loop without closing braces, its helper method, or the
\texttt{[Java End]} marker. Several vanilla compilation failures were direct
interface mismatches, including \texttt{iscube} instead of \texttt{isCube} in
HumanEval/77 and incompatible array/list parameter types in HumanEval/20, /62,
/120, /129, and /136. Supplying the exact Java declaration directly addressed
this pattern.

The crafted responses that compiled but failed tests illustrate the remaining
semantic difficulty. HumanEval/8 returned an array where the declaration
required \texttt{List<Integer>}; the resulting equality check did not have the
intended value semantics. HumanEval/145 applied the sign to the whole digit sum
instead of only the leading digit. HumanEval/129 reimplemented the description
using a global grid minimum rather than translating the neighbour-based loop,
and HumanEval/136 reversed the required extrema for negative and positive
values. HumanEval/148, in contrast, improved from a behavioral failure to a
pass. These examples suggest that source-code length alone is not decisive.
The persistent risk is re-solving from the description or overlooking a
Python-specific value, ordering, or state-transition rule instead of
translating the implementation statement by statement.

\subsection{MP3-Task2: Bug Detection}

\subsubsection{Evaluation and Success Rate}\hfill\break
For each HumanEvalPack program, the model reviewed the buggy implementation and
predicted either \texttt{Buggy} or \texttt{Correct}. Because every selected
HumanEvalPack entry contains an injected bug, a tagged \texttt{Buggy} verdict is
correct. The parser accepted one unambiguous verdict enclosed by
\texttt{<start>} and \texttt{<end>}. Vanilla prompting detected 2 of 20 bugs
(10\%), whereas the crafted prompt detected 15 of 20 (75\%), an absolute gain
of 65 pp. This exceeds the expected range, so we report both the semantic and
format-compliance effects below. Table~\ref{tab:mp3t2pairwise} gives the
pairwise results.

\begin{center}
\begin{minipage}{\columnwidth}\centering
\small
\setlength{\tabcolsep}{4pt}
\begin{tabular}{@{}lcc@{}}
\toprule
\texttt{problem\_ID} & \texttt{vanilla\_is\_correct} & \texttt{crafted\_is\_correct} \\
\midrule
HumanEval/60  & False & \textbf{True} \\
HumanEval/24  & True  & True          \\
HumanEval/67  & False & \textbf{True} \\
HumanEval/145 & False & False         \\
HumanEval/129 & False & \textbf{True} \\
HumanEval/77  & False & False         \\
HumanEval/79  & False & \textbf{True} \\
HumanEval/20  & False & \textbf{True} \\
HumanEval/62  & False & False         \\
HumanEval/102 & False & \textbf{True} \\
HumanEval/8   & False & False         \\
HumanEval/81  & False & \textbf{True} \\
HumanEval/140 & False & False         \\
HumanEval/161 & False & \textbf{True} \\
HumanEval/136 & False & \textbf{True} \\
HumanEval/75  & True  & True          \\
HumanEval/148 & False & \textbf{True} \\
HumanEval/118 & False & \textbf{True} \\
HumanEval/126 & False & \textbf{True} \\
HumanEval/120 & False & \textbf{True} \\
\midrule
Correct       & 2/20  & 15/20         \\
Success rate  & 10\%  & \textbf{75\%} \\
\bottomrule
\end{tabular}
\par\smallskip
\refstepcounter{table}\label{tab:mp3t2pairwise}\textbf{Table \thetable:}
Pairwise bug-detection results. Bold marks a False$\rightarrow$True change.
\end{minipage}
\end{center}

The two bugs detected by the vanilla prompt, HumanEval/24 and /75, were also
detected by the crafted prompt. Thirteen additional programs changed from
incorrect to correct, five remained incorrect, and none regressed. Unlike Task
1, where one crafted translation lost a previous success, the Task 2
improvement was monotonic over this dataset.

\subsubsection{Prompt Changes}\hfill\break
The vanilla prompt supplied the implementation, requested step-by-step
reasoning, and asked for a tagged verdict. The crafted prompt made four main
changes. First, it included the natural-language specification and behavioral
tests instead of showing only the implementation. Second, it supplied one
buggy and one correct demonstration, which reduced bias toward either label.
Third, it required exact comparison of operators, loop bounds, branches, and
return values, followed by manual evaluation of assertions and checks for
boundaries, negative-number behavior, mutation, ordering, and duplicates.
Fourth, it placed exactly one tagged verdict on the first line and limited the
remaining explanation to 250 words.

Including the behavioral tests is consistent with the task definition: the
correctness specification is represented by both the natural-language
description and its assertions. The tests therefore act as specification
context for reasoning, not as an external execution result or an answer key;
the model must still evaluate an assertion to discover a mismatch.

The format constraint was important. In the vanilla run, 14 responses were
parse errors, 4 explicitly predicted \texttt{Correct}, and only 2 correctly
predicted \texttt{Buggy}. In the final crafted run there were no parse errors:
15 responses correctly predicted \texttt{Buggy}, while 5 predicted
\texttt{Correct}. Therefore, the 65 pp gain was produced by both more reliable
verdict formatting and better use of the specification and tests.

These changes were applied as one prompt configuration, so the experiment does
not isolate the contribution of each individual change. In particular, the
behavioral tests provide concrete input-output obligations, while the
verdict-first format affects whether a conclusion can be parsed. We therefore
interpret the 65 pp increase as the effect of the complete crafted prompt, not
as the effect of the demonstrations or formatting instruction alone.

The gain is not explained only by strict parsing. Reading the 14 untagged
vanilla responses leniently, 13 still concluded that the implementation was
correct, and only HumanEval/129 described it as buggy. The vanilla success rate
would therefore be only 10--15\% even under a prose-based interpretation,
compared with 75\% for the crafted prompt. The improvement exceeds 50 pp
because the vanilla baseline left 90 pp of available headroom, not because the
parser discarded a large set of semantically correct bug predictions.

\subsubsection{When Bug Detection Succeeds or Fails}\hfill\break
All 13 vanilla-to-crafted improvements occurred without a regression. The
crafted approach worked best when a supplied assertion exposed a concrete
mismatch that could be found by tracing exact operators, bounds, or returned
values. The remaining misses were HumanEval/8, /62, /77, /140, and /145. They
span aggregation, polynomial-derivative logic, numerical cube checking,
character-by-character spacing, and digit-sum sorting. This diversity suggests
that program length alone is not the main difficulty. The model still struggles
when detecting the injected defect requires a carefully chosen counterexample
or precise reasoning about numerical and ordering edge cases. In all five final
misses the response produced an explicit \texttt{Correct} verdict, so these
were semantic misses rather than parser failures.

At the mutation level, the five misses comprised two value-misuse bugs,
HumanEval/62 (omitting the \texttt{[1:]} slice) and /8 (initializing
\texttt{prod\_value} to zero); two excess-logic bugs, /145 (an incorrect sign
expression) and /140 (an extra \texttt{"\_\_"} case); and one missing-logic
bug, /77 (omitting \texttt{** 3}). Each is a small edit, and the docstring
examples supplied to the model expose the incorrect behavior. However, bug
type alone did not predict success because the 15 detected bugs also span the
dataset's value, excess-logic, and missing-logic categories.

The crafted prompt was most effective when a single supplied assertion could
serve as a counterexample: after finding one mismatch, the model could stop and
label the implementation buggy. It was less reliable when the defect depended
on an unusual numeric value, a tie or ordering rule, or a state that emerges
only after several iterations. Those cases require the model not merely to
paraphrase the implementation, but to execute it precisely and compare the
observed value with the specification.

\subsubsection{Reasoning Versus Final Verdict}\hfill\break
Inspection of the final responses revealed discrepancies in both directions.
First, a correct \texttt{Buggy} label did not always mean that the explanation
identified the injected defect. Of the 15 crafted \texttt{Buggy} verdicts,
HumanEval/60, /24, and /136 clearly described the relevant mutation, and /81
described it more loosely. Many other responses reached the right label for an
unrelated or generic reason. For example, HumanEval/148 blamed planet ordering
rather than the \texttt{"Jupyter"} typo; /75 claimed that the product was never
compared with the input rather than identifying the changed
\texttt{range(0,n)} in \texttt{is\_prime}; /102 stated that 28 is not even; and
/126 did not discuss the removed duplicate check. These are correct
predictions without reasoning tied to the actual specification violation.

Second, all five crafted \texttt{Correct} predictions claimed that the
behavioral tests passed without evaluating a concrete assertion, despite the
prompt's instruction to do so. HumanEval/77 restated the buggy comparison in
prose and still concluded \texttt{Correct}. HumanEval/8 said that the function
returns a product of one for an empty list even though the implementation
initializes \texttt{prod\_value} to zero. In these cases, the explanation
followed the intended specification rather than the code that was actually
under review.

The final crafted run had no parser failures, but not every response followed
the requested presentation. HumanEval/161 placed its verdict at the end rather
than on the first line, and /81 embedded the tag inside a sentence. Several
responses also exceeded the 250-word limit. They were still parseable because
the evaluator requires one unambiguous tag anywhere in the response; it does
not enforce tag position or explanation length. Thus, ``no parse errors'' is
not equivalent to complete instruction compliance.

The vanilla responses show why strict parsing remains necessary. HumanEval/129
called the implementation buggy and proposed a correction but emitted no tag,
so the prediction could not be recorded unambiguously. HumanEval/60 reasoned
that \texttt{sum\_to\_n(100)} sums only through 99 but then gave an inconsistent
value and concluded that the function was correct. HumanEval/140 printed the
literal phrase \texttt{<start> and <end>} in addition to its verdict. We
therefore retained the single-unambiguous-tag rule and kept format failures
separate from semantic failures instead of inferring labels from free-form
prose.

\subsection{Lessons Across the Two MP3 Tasks}

The two tasks benefited from different kinds of added context. Translation
improved mainly when the crafted prompt fixed the target-language interface and
highlighted cross-language semantics. Bug detection improved when the prompt
made the specification executable through concrete assertions and made the
output contract unambiguous. In both tasks, structured output was necessary
for automated evaluation but did not guarantee semantic correctness.

The failure categories also reveal different limits. Task 1 uses compilation
and execution as external checks, so an answer can fail before any test runs or
can compile and still behave incorrectly. Task 2 evaluates the model's own
verdict, so a confidently formatted response may still overlook the injected
bug. Across both tasks, the strongest prompt changes supplied task-specific
constraints and a concrete checking procedure rather than merely requesting
more explanation.

\ifmpthreestandalone
  \end{document}
\fi
